% Project-specific proofs, included verbatim in each Towers report.
\section{Proofs for the multiple-test model and its numerical interpretation}
\label{sec:towers-proofs}

This appendix proves the deductive claims of the Towers study. Its probability
space is the ideal model, rather than a finite pseudorandom generator or an
empirical population. The foundational results and their proofs are in
Appendix~\ref{sec:proof-foundations}; the sampling results and their proofs are
in Appendix~\ref{sec:shared-statistical-methods}. The underlying standard
results are change of variables, Tonelli's theorem, dominated convergence,
the fundamental theorem of calculus and the intermediate value theorem
\citep{durrett2019,gallager2013}. Numerical tables are approximate evaluations
of the exact expressions below. Source dates, source quotations, software
inventories and retained execution counts are records to inspect, rather
than conclusions of probability theory.

\begin{proposition}[The Gaussian construction, density and joint tail]
\label{tw:gaussian}
Let \(G,E\) be independent standard normal variables, let \(|\rho|<1\), and
put \(X=G\), \(Y=\rho G+\sqrt{1-\rho^2}E\). Then \(X,Y\) have zero means,
unit variances, correlation \(\rho\), and joint density
\begin{equation}
 f_\rho(x,y)=\frac{1}{2\pi\sqrt{1-\rho^2}}
 \exp\!\left[-\frac{x^2-2\rho xy+y^2}{2(1-\rho^2)}\right].
 \label{eq:tw-density}
\end{equation}
For finite \(a,b\), their joint exceedance probability is
\begin{equation}
 J(a,b;\rho)=\int_a^\infty\varphi(x)
 \overline\Phi\!\left(\frac{b-\rho x}{\sqrt{1-\rho^2}}\right)\,dx.
 \label{eq:tw-tail-proof}
\end{equation}
For \(0<p<1\), the unique cutoff with exceedance probability \(p\) is
\(\Phi^{-1}(1-p)\). Strict and non-strict versions have equal probability.
\end{proposition}
\begin{proof}
Linearity gives the two means. Independence gives \(E[GE]=E[G]E[E]=0\);
expansion of the squares gives
\(\operatorname{Var}(Y)=\rho^2+(1-\rho^2)=1\) and
\(\operatorname{Cov}(X,Y)=\rho\). Gaussian affine closure is proved in
Proposition~\ref{pf:gaussian}. Alternatively the invertible transformation has
inverse \(g=x\), \(e=(y-\rho x)/\sqrt{1-\rho^2}\) and Jacobian determinant
\((1-\rho^2)^{-1/2}\). Multiplying \(\varphi(g)\varphi(e)\) by that
determinant and expanding the exponent proves \eqref{eq:tw-density}.
Dividing it by the positive marginal density \(\varphi(x)\) gives the
conditional density of \(Y\) at fixed \(X=x\): it is normal with mean
\(\rho x\) and variance \(1-\rho^2\). This is the continuous conditional
version explained in Proposition~\ref{pf:conditioning}; an event \(X=x\) itself
has probability zero. Tonelli's theorem integrates its upper tail against
\(\varphi(x)\), proving \eqref{eq:tw-tail-proof}. Finally \(\Phi\) is
continuous, strictly increasing, and has limits zero and one, because
\(\varphi>0\) and integrates to one. Inverting it proves the cutoff claim.
Its continuous density assigns each cutoff point probability zero.
\end{proof}

\begin{proposition}[Every denominator and the binary correlation]
\label{tw:events}
Write \(A=\{X>a\}\), \(B=\{Y>b\}\), \(p_A=P(A)\), \(p_B=P(B)\),
and \(J=P(A\cap B)\). The four disjoint cell probabilities, in the order
neither, \(A\) only, \(B\) only, both, are
\begin{equation}
 1-p_A-p_B+J,\quad p_A-J,\quad p_B-J,\quad J.
 \label{eq:tw-cells}
\end{equation}
Thus \(u=P(A\cup B)=p_A+p_B-J\), and
\begin{align}
 P(B\mid A)&=J/p_A,&P(A\mid B)&=J/p_B,\\
 P(A^c\mid B)&=1-J/p_B,&
 P(A\cap B\mid A\cup B)&=J/u.
 \label{eq:tw-denominators}
\end{align}
Here all denominators are positive for finite Gaussian cutoffs.
The Pearson correlation of the indicators is
\begin{equation}
 \operatorname{Corr}(1_A,1_B)=
 \frac{J-p_Ap_B}{\sqrt{p_A(1-p_A)p_B(1-p_B)}}.
 \label{eq:tw-phi-proof}
\end{equation}
At equal margins \(p\), exactly one pass has probability \(2(p-J)\),
one specified pass without the other has probability \(p-J\), and the
both-pass fraction of the eligible union is \(J/(2p-J)\).
\end{proposition}
\begin{proof}
The four cells partition the whole space. Subtracting the both-pass cell
from each marginal and then subtracting the other three cells from one
proves \eqref{eq:tw-cells}. Their nonnegativity proves the Fr\'echet bounds
\(\max(0,p_A+p_B-1)\leq J\leq\min(p_A,p_B)\), and summing the three
passing cells proves the union formula. In every displayed conditional,
intersect the numerator event with the conditioning event and divide by
the probability of that conditioning event, as in Proposition~\ref{pf:events}.
Since \(1_A^2=1_A\), its mean and variance are \(p_A\) and
\(p_A(1-p_A)\); likewise for \(B\). Also \(1_A1_B=1_{A\cap B}\), so
the covariance is \(J-p_Ap_B\). The correlation definition now proves
\eqref{eq:tw-phi-proof}. Substitution of equal margins proves the last
three assertions. These calculations do not identify binary correlation
with \(\rho\); each is a correlation of different random variables.
\end{proof}

\begin{proposition}[Correlation differentiation and sensitivity]
\label{tw:plackett}
For finite cutoffs and \(-1<\rho<1\),
\begin{equation}
 \frac{\partial J(a,b;\rho)}{\partial\rho}=f_\rho(a,b)>0,
 \qquad
 J(a,b;\rho)=\overline\Phi(a)\overline\Phi(b)
                 +\int_0^\rho f_t(a,b)\,dt.
 \label{eq:tw-plackett-proof}
\end{equation}
Consequently joint passing and either conditional passing probability
strictly increase with \(\rho\) at fixed marginal cutoffs, whereas the
union probability strictly decreases. Equal cutoffs give equal theoretical
conditionals; mixed cutoffs need not.
\end{proposition}
\begin{proof}
This is the bivariate case of Plackett's identity \citep{plackett1954}; a
derivation is included. With \(D=1-t^2\), direct differentiation gives
\[
 \partial_t f_t(x,y)=
 \left\{\frac{t}{D}+\frac{(1+t^2)xy-t(x^2+y^2)}{D^2}\right\}f_t(x,y)
 =\partial_x\partial_y f_t(x,y).
\]
On every compact correlation interval \(|t|\leq r<1\), the derivative is
bounded by a constant times
\((1+x^2+y^2)\exp[-(x^2+y^2)/(2(1+r))]\), an integrable function.
Dominated convergence therefore permits differentiation of the double
tail integral. Integrating \(\partial_x\partial_y f_t\) over
\([a,\infty)\times[b,\infty)\) gives \(f_t(a,b)\): the Gaussian
density and its first derivatives vanish at the infinite boundaries.
At \(t=0\) the density factors, so the joint tail is the product of the
marginal tails. The fundamental theorem of calculus gives the second
identity, including negative \(\rho\) with its oriented integral.
The density is strictly positive. The monotonicity claims follow by
dividing by the fixed positive margins, or subtracting \(J\) from their
sum. Swapping \(x,y\) leaves \(f_\rho\) unchanged, which proves the
equal-cutoff symmetry. At mixed margins the conditional ratio is exactly
\((J/p_A)/(J/p_B)=p_B/p_A\), explaining the factor of twenty in the
top-2\% versus top-0.1\% example.
\end{proof}

\begin{proposition}[A distribution-free contradiction and a Gaussian bound]
\label{tw:original-bound}
Marginal pass rates \(0.02\) and \(0.001\) cannot have union probability
\(0.0211\). Moreover, equal top-2\% Gaussian cutoffs at \(\rho=0.7\)
satisfy
\begin{equation}
 J>\frac{47}{15000}>0.0031,
 \qquad P(B\mid A)>\frac{47}{300}>0.156,
 \qquad P(A\cup B)<\frac{553}{15000}<0.0369.
 \label{eq:tw-rigorous-rejection}
\end{equation}
These exact inequalities already reject the first letter's overlap
\(0.002\), repeat-pass \(0.1027\) and union \(0.0383\) as model values,
without relying on the last digits of a numerical integrator.
\end{proposition}
\begin{proof}
The union bound gives \(P(A\cup B)\leq0.021<0.0211\).
For the Gaussian claims write \(T(x)=\overline\Phi(x)\). Integration of
\(t\varphi(t)=-\varphi'(t)\) gives the strict Mills upper bound
\(T(x)<\varphi(x)/x\) for \(x>0\), since \(t/x>1\) for \(t>x\).
For completeness the constant bound \(\pi>25/8\) also has an elementary
proof. The exact finite geometric identity
\[
 \frac1{1+t^2}=1-t^2+t^4-t^6+\frac{t^8}{1+t^2}
\]
and integration on \([0,x]\), \(x>0\), give
\(\arctan x>x-x^3/3+x^5/5-x^7/7\), because the remainder integral is
positive. The tangent addition formula gives
\(\arctan(1/2)+\arctan(1/3)=\pi/4\): its tangent is one and the
sum is in \((0,\pi/2)\), since its cosine has the positive factor
\(1-(1/2)(1/3)\). Applying the lower bounds at \(1/2\) and \(1/3\)
therefore gives
\[
 \frac\pi4>
 \sum_{x\in\{1/2,1/3\}}\left(x-\frac{x^3}{3}+\frac{x^5}{5}-\frac{x^7}{7}\right)
 =\frac{25}{32}+\frac{7765}{1959552}>\frac{25}{32}.
\]
Together with
\(e^{8/3}>\sum_{n=0}^4(8/3)^n/n!=3035/243>11\), this implies
\[
 T(7/3)<\frac{\varphi(7/3)}{7/3}
 <\frac{(2/5)}{11(7/3)}=\frac6{385}<\frac1{50}.
\]
Thus the cutoff \(a\) with \(T(a)=1/50\) is below \(7/3\).
Taylor's theorem with its signed third derivative gives
\(e^{-v}\leq1-v+v^2/2\) for \(v\geq0\). Hence
\[
 T(1)=\frac12-\int_0^1\varphi(t)\,dt
 >\frac12-\frac25\left(1-\frac16+\frac1{40}\right)
 =\frac{47}{300}.
\]
For every \(x\geq a\), the argument of the conditional tail in
\eqref{eq:tw-tail-proof} is at most
\(0.3a/\sqrt{0.51}<0.7/\sqrt{0.51}<1\). Therefore
\(J>T(a)T(1)>47/15000\). Division by \(1/50\) and the equal-margin
union identity prove the remaining inequalities.
\end{proof}

\begin{proposition}[Same-score nesting and different-test non-nesting]
\label{tw:nesting}
If \(b>a\) are applied to the same score \(Z\), then
\(\{Z>b\}\subset\{Z>a\}\). When \(P(Z>b)>0\), the conditional
probability of passing the lower cutoff given the higher is one.
With different Gaussian scores
and \(|\rho|<1\), neither finite-cutoff pass event contains the other
up to null sets.
\end{proposition}
\begin{proof}
The first inclusion is the elementary implication \(z>b\Rightarrow z>a\).
Its conditional probability is therefore the numerator probability divided
by itself. For different scores the open rectangles
\((a,a+1)\times(b-2,b-1)\) and
\((a-2,a-1)\times(b,b+1)\) lie respectively in \(A\setminus B\) and
\(B\setminus A\). The density \eqref{eq:tw-density} is continuous and
positive on their closures, so each rectangle has positive probability.
This proves non-nesting, strict union enlargement over either margin, and
strictly subunit conditional probabilities. The statement uses finite
cutoffs and a nonsingular Gaussian distribution; it does not assert a
relationship between actual admission policies or actual tests.
\end{proof}

\begin{proposition}[Existence and uniqueness of OR calibration]
\label{tw:calibration}
For every \(-1<\rho<1\) and \(0<u_0<1\), there is exactly one finite
\(t\) satisfying
\(2\overline\Phi(t)-J(t,t;\rho)=u_0\).
On a reporting scale \(I=100+16X\), its cutoff is \(100+16t\).
\end{proposition}
\begin{proof}
Let \(U(t)=P(X>t\text{ or }Y>t)\). Inclusion-exclusion identifies it with
the left side. If \(t_1<t_2\), the event at \(t_2\) is contained in the
event at \(t_1\). A rectangle with \(t_1<X<t_2\) and \(Y<t_1\) has
positive probability and belongs to the set difference; consequently
\(U(t_1)>U(t_2)\). For \(t_n\to t\), the event indicators converge
except on \(X=t\) or \(Y=t\), which have probability zero. Dominated
convergence proves continuity. As \(t\to\infty\), the union bound gives
\(U(t)\leq2\overline\Phi(t)\to0\). As \(t\to-\infty\),
\(U(t)\geq\overline\Phi(t)\to1\). The intermediate value theorem gives
existence and strict decrease gives uniqueness. A positive affine reporting
scale transforms the threshold by the same affine formula; its mean and
standard deviation are 100 and 16 by Proposition~\ref{pf:moments}.
\end{proof}

\begin{proposition}[A rarity equivalent identifies a size, not a set]
\label{tw:rarity}
For \(0<p<1\), \(100+16\Phi^{-1}(1-p)\) is the unique SD-16 normal
cutoff exceeded with probability \(p\). Equal selected fractions do not
imply equal selected people or identify any latent ability.
Normal margins and their correlation do not identify the Gaussian overlap.
\end{proposition}
\begin{proof}
The cutoff claim follows by inversion in Proposition~\ref{tw:gaussian}
and the positive affine conversion. To disprove set identification, let
\(U\) be uniform on \([0,1)\). The intervals \([0,p)\) and the same
interval translated modulo one by \(d=\min(p,1-p)/2>0\) both have
probability \(p\), but their symmetric difference has positive length.
Thus equal probability does not identify the event. For latent-variable
identification, keep the same observed pair \((X,Y)\) and the same observed
selection event \(A=\{X>a\}\). In one model
adjoin an independent standard normal variable \(L\), called latent ability.
In another model set \(L=X\). Both have the identical observed score law
and a standard normal \(L\), but observed selection is independent of
\(L\) in the first model and associated with it in the second. No
measurement relation specifying \(L\) was part of the observed model, so
the observations cannot choose between them. Finally
Proposition~\ref{pf:normal-marginals} constructs normal margins with the same
correlation using \(Y=SX\) with an independent biased sign \(S\).
For positive cutoffs its joint tail is
\((1+\rho)\overline\Phi(\max(a,b))/2\), not generally
\eqref{eq:tw-tail-proof}. This provides the stated explicit failure of
identification by margins and correlation alone.
\end{proof}

\begin{proposition}[Repeated routes and selection]
\label{tw:attempts}
For any finite pass events \(A_1,\ldots,A_m\),
\[
 \max_iP(A_i)\leq P(\cup_iA_i)\leq\min(1,\sum_iP(A_i)).
\]
If they are independent with common probability \(p\), the union is
\(1-(1-p)^m\). A population correlation does not, by itself, determine a
correlation in a selected subpopulation.
\end{proposition}
\begin{proof}
Each individual event is contained in the union, proving the lower bound.
Pointwise \(1_{\cup_iA_i}\leq\sum_i1_{A_i}\); expectations and the
universal bound one prove the upper bound. Independence of the complements
gives \(P(\cap_iA_i^c)=(1-p)^m\), and taking the complement proves the
formula. For a concrete selection counterexample take independent standard
normals \(X,Y\), and condition on \(XY>0\). This event has probability
one half; symmetry gives zero conditional means. Conditional squares retain
mean one, but the conditional product has expectation
\(E[|X|]E[|Y|]=2/\pi>0\). Indeed the signs are independent of the
absolute values, and \(E|X|=2\int_0^\infty x\varphi(x)\,dx
=\sqrt{2/\pi}\). The selected correlation is consequently \(2/\pi\),
despite population correlation zero. This is an existence counterexample,
not a model of any named society's selection.
\end{proof}

\begin{proposition}[What the BASIC grid normalizes]
\label{tw:grid}
Let \(h=1/2\), let \(F=\{1,1+h,\ldots,199\}^2\), and give node
\((x,y)\) weight \(g(x,y)>0\) proportional to its Gaussian IQ-scale
density. The ratio of the sum over passing nodes to the sum over \(F\)
is exactly a probability for a finite weighted distribution on \(F\).
It is not exactly a continuous Gaussian tail. There are \(397^2=157609\)
nodes; constant density-scale factors and \(h^2\) cancel from the ratio.
Viewed as composite midpoint quadrature, nodes starting at \(Q\) cover
cells starting at \(Q-h/2=Q-0.25\), not at \(Q\).
\end{proposition}
\begin{proof}
The normalized node masses are
\(w_{xy}=g(x,y)/\sum_Fg\); they are positive and sum to one. Summing
them over the passing subset proves the discrete-probability assertion and
also supplies its own consistent marginal denominator. The node count in
each coordinate is \((199-1)/(1/2)+1=397\). Multiplying every weight by
any common positive constant leaves every normalized mass unchanged,
proving cancellation of the cell area and the density Jacobian factor
\(16^{-2}\). Centering a cell of width \(h\) on its node gives the
stated lower boundary and upper boundary \(199+h/2\).

More precisely, let \(N,D\) be the continuous Gaussian masses of the
passing cell rectangle and the full cell rectangle, and let \(n,d\) be
their midpoint sums including cell area. On a finite rectangle a twice
continuously differentiable density satisfies the bound
\[
 |n-N|\leq\frac{|R|h^2}{24}(M_{xx}+M_{yy}),
\]
where \(M_{xx},M_{yy}\) bound the absolute second derivatives on the
rectangle. Apply the midpoint remainder first in \(x\), then in \(y\):
the linear Taylor term integrates to zero and
\(\int_{-h/2}^{h/2}u^2/2\,du=h^3/24\); summing the cells proves the
bound. The same argument bounds \(|d-D|\). Since \(d>0\),
\[
 \left|\frac nd-\frac ND\right|
 \leq\frac{|n-N|+(N/D)|d-D|}{d}.
\]
The truncated continuous ratio \(N/D\) is itself a conditional probability
given the full finite cell rectangle. Its excluded mass is bounded by the
sum of the four marginal tails beyond its edges, by the union bound.
This separates discretization, shifted boundary and truncation; it does
not claim that their errors vanish or have the same sign at finite \(h\).
Comparisons of their sizes in the tables are numerical comparisons, not a
certified integration-error bound. An integer cutoff \(Q\) has actual
continuous marginal \(\overline\Phi((Q-100)/16)\); division by a
different nominal tail answers no conditional probability of this model.
Rounding the scores first would define yet another event and requires an
explicit rounding convention.
\end{proof}

\begin{proposition}[Quantization of a conditional ratio]
\label{tw:quantization}
If \(|j'-j|\leq\delta\), \(|p'-p|\leq\delta\), \(p>\delta\),
and \(j\geq0\), then
\[
 \left|\frac{j'}{p'}-\frac jp\right|
 \leq\frac{\delta(1+j/p)}{p-\delta}.
\]
If also \(0\leq j\leq p\), this is at most
\(2\delta/(p-\delta)\).
\end{proposition}
\begin{proof}
Put \(\Delta j=j'-j\), \(\Delta p=p'-p\). The difference is
\((p\Delta j-j\Delta p)/(pp')\), and \(p'\geq p-\delta>0\).
The triangle inequality bounds the numerator by \(p\delta+j\delta\).
Division gives the first bound, and \(j/p\leq1\) gives the second.
For \(\delta=0.51\times10^{-15}\) and \(p=1/33000\), substituting
in the latter expression gives less than \(3.37\times10^{-11}\).
This accounts for probability quantization only, not integration error.
\end{proof}

\begin{proposition}[Simulation partition and selected denominators]
\label{tw:simulation}
For \(T\) simulated pairs, let \(n_{rs}\) count the four indicator
values \((1_A,1_B)=(r,s)\). Then
\[
 \sum_{r,s}n_{rs}=T,\quad n_A=n_{10}+n_{11},\quad
 n_B=n_{01}+n_{11},\quad n_U=n_A+n_B-n_{11}.
\]
Under ideal independent identical draws, each fixed event count is binomial.
Given \(n_A=m>0\), \(n_{11}\) is binomial with \(m\) trials and
success probability \(J/p_A\). Different cutoff rules evaluated on the
same pairs need not give independent estimates.
\end{proposition}
\begin{proof}
For each pair exactly one of the four Boolean outcomes occurs. Summing that
pointwise partition proves the integer identities without any distributional
assumption. Under the additional independent-pair model, a fixed event gives
independent Bernoulli indicators, whose sum has the binomial law by
Proposition~\ref{sm:binomial}. Conditional on the identities of the \(m\)
observations passing \(A\), their \(B\) indicators are independent with
conditional probability \(J/p_A\). Every choice of these identities gives
the same binomial sum; mixing over the choices proves the claim given
\(n_A=m\), as proved generally in Proposition~\ref{sm:conditional-count}.
When \(m=0\), the conditional frequency is undefined. To see dependence
between common-pair rules explicitly, if one pass event \(C\) is contained
in another \(D\) and both have probabilities strictly between zero and
one, their one-pair indicators have covariance
\(P(C)(1-P(D))>0\). Their frequency estimates from \(T\) independent
pairs have this covariance divided by \(T\), by expanding the double
sum and canceling distinct-pair covariances. Thus independence between rows
cannot be inferred from independent observations within each row.
\end{proof}
